%% sections/taxonomy.tex
\section{Three-Axis Taxonomy of 3DGS Research}
\label{sec:taxonomy}

We organize 3DGS research along three orthogonal axes that capture
the essential dimensions of the field:
\textbf{Task Domain} (what problem is being solved),
\textbf{Module Focus} (which component is improved), and
\textbf{Technique} (how the improvement is achieved).
Figure~\ref{fig:taxonomy} provides an overview.

\begin{figure*}[t]
    \centering
    % \includegraphics[width=\textwidth]{taxonomy.pdf}
    \todo{Insert taxonomy figure}
    \caption{Three-axis taxonomy of 3DGS research. Each axis represents
    a distinct categorization dimension, enabling multi-perspective
    analysis of the field.}
    \label{fig:taxonomy}
\end{figure*}

% ---- Axis 1: Task Domain ----
\subsection{Axis 1: Task Domain}
\label{sec:taxonomy:task}

3DGS methods can be categorized by the primary task they address:

\begin{itemize}[leftmargin=*]
    \item \textbf{Static Reconstruction}: The foundational task of
          reconstructing static 3D scenes from multi-view images
          \cite{kerbl3Dgaussians, yu2024mipsplatting}.
    \item \textbf{Dynamic / 4D}: Extending 3DGS to dynamic scenes with
          temporal variation \cite{wu20244dgs, yang2024deformable}.
    \item \textbf{Editing}: Modifying 3DGS scenes for appearance,
          geometry, or composition changes.
    \item \textbf{Generation}: Generating 3DGS scenes from text, images,
          or other modalities.
    \item \textbf{Compression}: Reducing the memory footprint of 3DGS
          representations for storage and transmission.
\end{itemize}

% ---- Axis 2: Module Focus ----
\subsection{Axis 2: Module Focus}
\label{sec:taxonomy:module}

The 3DGS pipeline consists of several modules, and improvements typically
target specific components:

\begin{itemize}[leftmargin=*]
    \item \textbf{Representation}: Modifying the Gaussian primitive
          structure (e.g., 2D Gaussians \cite{huang20242dgs},
          anchored Gaussians).
    \item \textbf{Optimization}: Improving the training loss, scheduling,
          or regularization.
    \item \textbf{Rendering}: Enhancing the rasterizer (e.g.,
          anti-aliasing \cite{yu2024mipsplatting}).
    \item \textbf{Initialization}: Better starting points for Gaussian
          positions and properties.
    \item \textbf{Densification}: Adaptive control of Gaussian count
          during training.
\end{itemize}

% ---- Axis 3: Technique ----
\subsection{Axis 3: Technique}
\label{sec:taxonomy:technique}

The specific techniques employed to improve 3DGS include:

\begin{itemize}[leftmargin=*]
    \item \textbf{Anchoring}: Using anchor points to structure Gaussian
          distribution \cite{yu2024scaffold}.
    \item \textbf{Compression}: Pruning, quantization, or hash-based
          encoding.
    \item \textbf{Anti-Aliasing}: Multi-scale filtering to prevent
          aliasing artifacts \cite{yu2024mipsplatting}.
    \item \textbf{Regularization}: Imposing geometric or photometric
          priors.
    \item \textbf{Level-of-Detail (LoD)}: Hierarchical representation
          for scalable rendering.
    \item \textbf{Feature Augmentation}: Adding semantic or geometric
          features to Gaussians.
\end{itemize}

\todo{Add cross-axis analysis table showing how techniques map to
modules and tasks.}
